\subsection{Definition of a vectorial measure}

The definition of a measure given in Ch. III, §1, No.~3 may be generalized as follows:

DEFINITION 1. --- Let F be a Hausdorff locally convex space over R. One calls vectorial measure on T with values in F every continuous linear mapping of the space \(\mathcal{H}(T)\) into F.

Def. 1 may also be expressed as follows: a vectorial measure on T with values in F is a linear mapping m of \(\mathcal{K}(T)\) into F such that, for every compact subset K of T, the restriction of m to \(\mathcal{K}(T,K)\) is continuous for the topology of uniform convergence. If \(f \in \mathcal{K}(T)\) , one also writes \(\int f dm\) or \(\int f(t) dm(t)\) instead of \(\mathbf{m}(f)\) . The measures with values in R are sometimes called real measures (Ch. III, §1, No.~5) or scalar measures on T.

Examples. --- 1) The identity mapping of \(\mathcal{H}(\mathrm{T})\) is a vectorial measure on \(\mathrm{T}\) with values in \(\mathcal{H}(\mathrm{T})\) .

\begin{enumerate}
\def\labelenumi{\arabic{enumi})}
\setcounter{enumi}{1}
\tightlist
\item
  *Let H be a complex Hilbert space, \(\mathcal{L}(\mathrm{H})\) the normed algebra of continuous endomorphisms of H. Let A be a subalgebra of \(\mathcal{L}(\mathrm{H})\) , commutative, closed, self-adjoint and containing the identity; one shows that there exist a compact space X and an isomorphism of the normed algebra A onto the algebra \(\mathcal{K}_{\mathbf{C}}(\mathrm{X}) = \mathcal{C}_{\mathbf{C}}(\mathrm{X})\) of continuous complex functions on X, equipped with the norm \(\| f \| = \sup_{x \in \mathrm{X}} |f(x)|\) . The inverse isomorphism, restricted to \(\mathcal{K}(\mathrm{X})\) , is a vectorial measure \(\mathbf{m}\) on X, with values in \(\mathcal{L}(\mathrm{H})\) , such that \(\mathbf{m}(fg) = \mathbf{m}(f)\mathbf{m}(g)\) .*
\end{enumerate}

Remarks. --- 1) For a linear mapping m of \(\mathcal{K}(\mathrm{T})\) into F to be a vectorial measure, it is necessary and sufficient that, for every compact subset K of T, the image under m of the unit ball \(\|f\|\leqslant1\) of \(\mathcal{K}(\mathrm{T},\mathrm{K})\) be bounded in F. The notion of vectorial measure with values in F is therefore the same for all the Hausdorff locally convex topologies on F that admit the same bounded sets, and in particular for all the topologies compatible with the duality between F and F' (TVS, IV, §1, No.~1, Prop. 1).

\begin{enumerate}
\def\labelenumi{\arabic{enumi})}
\setcounter{enumi}{1}
\item
  Let \(T_{1}\) be a locally compact space, \(F_{1}\) a Hausdorff locally convex space over R, u a continuous linear mapping of \(\mathcal{K}(T_{1})\) into \(\mathcal{K}(T)\) , and v a continuous linear mapping of F into \(F_{1}\) . If m is a vectorial measure on T with values in F, then \(v \circ m \circ u\) is a vectorial measure on \(T_{1}\) with values in \(F_{1}\) . In particular, if g is a continuous, finite numerical function on T, then \(f \mapsto \mathbf{m}(gf)\) is a vectorial measure on T with values in F, which is denoted \(g \cdot m\) ; if h is a second continuous, finite numerical function on T, then \(g \cdot (h \cdot \mathbf{m}) = (gh) \cdot \mathbf{m}\) .
\item
  Since the space \(\mathcal{H}(\mathrm{T})\) is the direct limit of the Banach spaces \(\mathcal{H}(\mathrm{T},\mathrm{K})\) , and is in particular barreled (TVS, III, §4, No.~1, Cor. of Prop. 2 and Cor. 3 of Prop. 3), in order that a linear mapping \(\mathbf{m}\) of \(\mathcal{H}(\mathrm{T})\) into \(\mathbf{F}\) be a vectorial measure, it is necessary and sufficient that, for every \(\mathbf{z}' \in \mathbf{F}'\) , \(\mathbf{z}' \circ \mathbf{m}\) be a scalar measure on \(\mathrm{T}\) (TVS, II, §6, No.~4, Prop. 5 and IV, §1, No.~3, Prop. 7).
\item
  In view of Remark 1, Prop. 1 of Ch. III, §2, No.~1 and its proof are again valid for vectorial measures. One can therefore define the support of a vectorial measure \(\mathbf{m}\) on \(T\) to be the complement of the largest open set \(U \subset T\) such that the restriction of \(\mathbf{m}\) to \(U\) is zero.
\end{enumerate}

